% Vector redraw of the author's original Figure 1.
% The figure is intentionally kept as TikZ so that all labels are typeset
% by LaTeX and remain vector objects in the compiled manuscript PDF.
\begin{tikzpicture}[x=1cm,y=1cm,>=Stealth,line cap=round,line join=round]
\tikzset{
  panel/.style={draw=black,line width=.6pt,rounded corners=6pt},
  interface/.style={draw=black,line width=.7pt},
  axis/.style={draw=black,line width=.45pt},
  crack/.style={draw=black,double=white,double distance=2.0pt,line width=.45pt},
  processzone/.style={draw=black,line width=2.0pt},
  alternativezone/.style={draw=black,line width=1.2pt,dashed}
}

% Common panel geometry.  The interface rays are symmetric about the
% horizontal bisector of material 2, as in the author's original sketch.
\newcommand{\failurepanel}[2]{%
  \begin{scope}[xshift=#1cm]
    \draw[panel] (-2.20,-1.55) rectangle (2.20,1.55);
    \coordinate (O) at (0,0);
    \coordinate (Iu) at ($(O)+(118:2.05)$);
    \coordinate (Il) at ($(O)+(242:2.05)$);
    \draw[interface] (O)--(Iu);
    \draw[interface] (O)--(Il);
    \draw[axis] (-2.05,0)--(2.05,0);

    % Polar coordinates and the half-angle alpha.
    \draw[->,line width=.65pt] (O)--++(27:1.62)
      node[above right=-1pt,font=\small] {$r$};
    \draw[->,line width=.5pt] (0.58,0)
      arc[start angle=0,end angle=27,radius=.58];
    \node[font=\small,fill=white,inner sep=.7pt] at (0.93,0.18) {$\theta$};
    \draw[->,line width=.5pt] (0.83,0)
      arc[start angle=0,end angle=118,radius=.83];
    \node[font=\small,fill=white,inner sep=.7pt] at (0.16,0.88) {$\alpha$};

    \node[font=\small,fill=white,inner sep=.7pt] at (0.18,-0.20) {$O$};
    \node[font=\small] at (-1.45,-0.92) {$E_1,\,\nu_1$};
    \node[font=\small] at ( 1.30,-0.92) {$E_2,\,\nu_2$};
    #2
  \end{scope}%
}

% (a) Intact broken interface.
\failurepanel{0}{%
  \node[font=\small] at (0,-1.88) {(a)};
}

% (b) Symmetric interfacial shear cracks of length l.
\failurepanel{5.00}{%
  \draw[crack] (O)--++(118:1.30);
  \draw[crack] (O)--++(242:1.30);
  \node[font=\small] at (-0.77, 1.18) {$l$};
  \node[font=\small] at (-0.77,-1.18) {$l$};
  \node[font=\small] at (0,-1.88) {(b)};
}

% (c) Mode-I process-zone alternatives along the two material bisectors.
\failurepanel{10.00}{%
  \draw[crack] (O)--++(118:1.30);
  \draw[crack] (O)--++(242:1.30);
  \node[font=\small] at (-0.77, 1.18) {$l$};
  \node[font=\small] at (-0.77,-1.18) {$l$};
  \draw[processzone] (O)--(-1.35,0);
  \draw[alternativezone] (O)--(1.35,0);
  \node[font=\small,below=2pt] at (-0.75,0) {$d_1$};
  \node[font=\small,below=2pt] at ( 0.75,0) {$d_2$};
  \node[font=\small] at (0,-1.88) {(c)};
}
\end{tikzpicture}
